\section{What Does the Evidence Say About Founder Wealth?}

\subsection{Empirical Realities of Self-Employment and Equity}

The empirical literature is less romantic than the public conversation about startups.

Barton Hamilton's widely cited analysis of U.S. self-employment, published in 2000, found that median self-employed earnings after 10 years in business were substantially below comparable paid-employment earnings (approximately 35\% lower in his reported comparison).\footnote{Barton H. Hamilton, ``Does Entrepreneurship Pay? An Empirical Analysis of the Returns to Self-Employment,'' \textit{Journal of Political Economy} 108(3), 2000.} This does not prove that founding a modern software company is a bad investment, or that every self-employed person would earn more in a job. It does challenge the idea that becoming one's own boss normally creates superior monetary income.

Moskowitz and Vissing-Jørgensen, studying entrepreneurial investments in the United States, found that private entrepreneurial equity had historically offered returns that did not appear higher than public equity returns despite far poorer diversification. Their 2002 paper framed this as a puzzle: why do people hold so much concentrated private-business risk without an obvious average return premium?\footnote{Tobias J. Moskowitz and Annette Vissing-J\o{}rgensen, ``The Returns to Entrepreneurial Investment: A Private Equity Premium Puzzle?'' \textit{American Economic Review} 92(4), 2002.} Possible explanations include nonfinancial benefits and preferences that ordinary portfolio models do not capture. Their evidence has historical and measurement limitations, but the central warning remains relevant: ownership alone does not guarantee an attractive risk-adjusted return.

Hall and Woodward examined venture-capital-backed founders and emphasized the huge risk borne by entrepreneurs whose compensation depends on uncertain exits. Their model found that the high average payoff in a selected venture-backed sample could coexist with a much less attractive choice for a risk-averse person without substantial initial wealth.\footnote{Robert E. Hall and Susan E. Woodward, ``The Burden of the Nondiversifiable Risk of Entrepreneurship,'' \textit{American Economic Review} 100(3), 2010.} The mean payoff is not the experience of the typical founder.

Official U.S. establishment statistics provide another useful caution. The Bureau of Labor Statistics reported 5-year survival rates in the approximate 50\% to 57\% range across selected cohorts of new business establishments, with the 2018 cohort at 57.3\%.\footnote{U.S. Bureau of Labor Statistics, ``Business Employment Dynamics Twentieth Anniversary!'' (2024).} These are establishment-survival statistics, not probabilities of achieving venture success. Surviving does not mean earning a good return; closing does not necessarily mean losing everything; and U.S. historical rates should not be transplanted into an Indian business forecast. Still, the numbers dispel the assumption that a newly opened business is likely to remain operating indefinitely.

There is equally important evidence against the idea that exceptional company-building is primarily a young person's game. In a study of U.S. growth-oriented startups, Azoulay and coauthors found that the mean age of founders of the fastest-growing 0.1\% of new ventures was 45, and that relevant industry experience predicted success.\footnote{Pierre Azoulay, Benjamin Jones, J. Daniel Kim, and Javier Miranda, ``Age and High-Growth Entrepreneurship,'' \textit{American Economic Review: Insights} 2(1), 2020.} The lesson is not to wait until 45. It is that accumulated experience and judgment can be valuable entrepreneurial assets; a founder does not have to abandon a successful career prematurely to retain a chance of building a significant company.

These studies examine different populations: self-employed individuals, private business owners, venture-backed founders, establishments, and high-growth firms. Their numbers should not be combined into a single fictional statistic such as ``the probability that a startup succeeds.'' There is no universal number. A freelance architect, a bootstrapped software business, a drug-discovery startup, and a neighborhood restaurant have different outcome distributions.

\subsection{The 4 Grounded Takeaways}

What can we say confidently? 4 things:
\begin{enumerate}[label=\textbf{\arabic*.},leftmargin=11pt,itemsep=1.5pt,topsep=1.5pt,parsep=0pt]
  \item \textbf{Spectacular wealth is possible but atypical.}
  \item \textbf{Ordinary self-employment does not reliably outperform comparable employment in financial terms.}
  \item \textbf{Founder risk is concentrated and is not fully represented by a company's valuation.}
  \item \textbf{Better selection of problems, stronger capabilities, and disciplined execution can matter even though luck remains irreducible.}
\end{enumerate}

\textbf{The evidence should make a founder more careful, not less ambitious.}
